%%
%% UTF-8       UTF-8      UTF-8       UTF-8       UTF-8       UTF-8
%% UTF-8       UTF-8      UTF-8       UTF-8       UTF-8       UTF-8
%%
%% XeLaTeX  XeLaTeX  XeLaTeX  XeLaTeX  XeLaTeX  XeLaTeX  XeLaTeX  XeLaTeX  

% !TeX program = xelatex
\lccode`\-=`\-
\documentclass[12pt,a4paper,oneside]{extarticle}
\usepackage{xunicode} % https://ctan.org/pkg/xunicode
\usepackage{xltxtra} % https://ctan.org/pkg/xltxtra
\usepackage{fontspec} % https://ctan.org/pkg/fontspec
\usepackage{polyglossia} % https://ctan.org/pkg/polyglossia
\setdefaultlanguage{portuges}
\usepackage{csquotes}

% Compilar com XeLaTeX
\usepackage{fontspec}
%\setmainfont{LMRoman10-Regular}
%\setsansfont{LMSans10-Regular}
%\setmonofont{LMMono10-Regular}

\clubpenalty=100
\linepenalty=100
\widowpenalty=100


\usepackage[backend=biber, style=numeric]{biblatex} 
\addbibresource{bibliografia.bib}

%estética
% \usepackage[hyphens]{url}
\usepackage{xurl}
\urlstyle{same} % URLs na mesma fonte do texto (Times New Roman)
\AtBeginDocument{%
  \setlength{\Urlmuskip}{0mu plus 1mu} % Ajusta espaçamento entre caracteres em URLs
  \setlength{\emergencystretch}{3em} % Permite maior flexibilidade na justificação
  \def\UrlBreakPenalty{100} % Penalidade baixa para facilitar quebras
  \def\UrlBigBreakPenalty{50} % Penalidade para quebras em caracteres como /
  \sloppy % Relaxa a justificação para evitar ultrapassagem
}

\usepackage[hang,small,bf]{caption}
\usepackage[a4paper,left=2.5cm,right=2.5cm,top=2.5cm,bottom=2.5cm,includehead,includefoot]{geometry}
\linespread{1.3}

\usepackage[
  xetex,                      % Driver para XeLaTeX
  unicode=true,               % Suporte a caracteres Unicode (ex.: ç, ã)
  bookmarks=true,             % Cria bookmarks no PDF
  bookmarksopen=true,         % Bookmarks abertos por padrão
  bookmarksnumbered=true,     % Numera os bookmarks (ex.: 1, 1.1, 1.2)
  linktoc=all,                % Links clicáveis em sumário, figuras e citações
  breaklinks=true,            % Permite quebra de links longos
  colorlinks=true,            % Links com cor em vez de caixas
  citecolor=black,            % Cor das citações (preto)
  filecolor=black,            % Cor dos links de arquivos (preto)
  linkcolor=black,            % Cor dos links internos (preto)
  urlcolor=blue,              % URLs em azul
  pdfpagelabels=true,         % Adiciona rótulos às páginas do PDF
  plainpages=false,           % Diferencia numeração de páginas
  pdftitle={Instalação do Arch Linux com XFCE no Android},
  pdfauthor={Luís A.V. Ferreira}, % Nome do autor
  pdfsubject={Instalação do Arch Linux com XFCE no Android},
  pdfkeywords={Arch Linux XFCE Android},
  pdfcreator={XeLaTeX com hyperref},
  pdfproducer={XeLaTeX}
]{hyperref}

%% Definir um novo estilo de url
\makeatletter
\def\url@luisstyle{%
  \@ifundefined{selectfont}{\def\UrlFont{\sf}}{\def\UrlFont{\small\ttfamily}}}
\makeatother
\urlstyle{luis}

\setlength{\captionmargin}{1.5in}
\usepackage{fancyhdr}
%\setlength{\headheight}{1in}

\renewcommand{\thefootnote}{\alph{footnote}} %modificar as letras de notas de rodapé%
\makeatletter
\renewcommand\@makefnmark{\@textsuperscript{\normalfont(\@thefnmark)}}
\makeatother
\renewcommand{\fboxrule}{2pt} %para alterar a expessura das caixas
%\renewcommand{\fboxsep}{3ex} %para alterar a separação nas caixas

\long\def\symbolfootnote[#1]#2{\begingroup\def\thefootnote{\fnsymbol{footnote}}\footnote[#1]{#2}\endgroup}

%%%%%%%%%%%%%%%%%%%%%%%%%%%%
%%%%%%%%%%%%%%%%%%%%%%%%%%%%
%%%%%%%%%%%%%%%%%%%%%%%%%%%%
\usepackage{setspace}
%\singlespacing
%\onehalfspacing
%\doublespacing
\usepackage{comment}
\usepackage{siunitx}
\usepackage{lettrine}
\usepackage{GoudyIn}
%%%%%%%%%%%%%%%%%%%%%%%%%%%%
\usepackage[x11names]{xcolor} 
\renewcommand{\LettrineFontHook}{\color{VioletRed4}\GoudyInfamily{}}
\LettrineTextFont{\itshape}
\setcounter{DefaultLines}{3}%
%%%%%%%%%%%%%%%%%%%%%%%%%%%%
\usepackage{relsize}
%\newcommand\CC{C\nolinebreak\hspace{-.05em}\raisebox{.4ex}{\relsize{-3}{\textbf{+}}}\nolinebreak\hspace{-.10em}\raisebox{.4ex}{\relsize{-3}{\textbf{+}}}}
\newcommand\CC{C\nolinebreak[4]\hspace{-.05em}\raisebox{.4ex}{\relsize{-3}{\textbf{++}}}}

%\usepackage{showframe}
\usepackage{pgf,tikz}       % - Elementos gráficos
\usepackage{tikzpagenodes}
\usetikzlibrary{backgrounds,calc,positioning}
\usetikzlibrary{positioning, shapes}

\pgfdeclarelayer{background}
\pgfdeclarelayer{foreground}
\pgfsetlayers{background,main,foreground}

\addbibresource{bibliografia.bib}

% \addbibresource{biblio.bib}
% \addbibresource{legislacao.bib}


\usepackage{listings}
\usepackage{lstlinebgrd}

% Definir cores personalizadas
\definecolor{lightgray}{gray}{0.95}
\definecolor{darkgray}{gray}{0.4}
\definecolor{purple}{rgb}{0.58,0,0.82}

% Traduzir o título "Listing" -> "Listagem"
\renewcommand{\lstlistingname}{Listagem}
\renewcommand{\lstlistlistingname}{Lista de Listagens}

% Definição do estilo
\lstset{
  basicstyle=\ttfamily\small,
  keywordstyle=\color{blue},
  stringstyle=\color{red!70!black},
  commentstyle=\color{green!50!black},
  breaklines=true,
  showstringspaces=false,
  upquote=true,
  extendedchars=true,
  literate={~}{{\textasciitilde}}1,
  numbers=left,                     % numeração das linhas à esquerda
  numberstyle=\tiny\color{darkgray},% estilo dos números
  numbersep=8pt,                    % distância entre número e código
  backgroundcolor=\color{lightgray},% cor de fundo base
  frame=single,                     % moldura à volta do código
  rulecolor=\color{black!30},       % cor da moldura
  language=python,                  % linguagem predefinida
  captionpos=t,                     % legenda em baixo
  xleftmargin=2em,                  % margens adicionais para espaço
  aboveskip=1em, belowskip=1em,
  linebackgroundcolor={%
        \ifodd\value{lstnumber}%
            \color{LightYellow1}%
        \else%
            \color{DarkSeaGreen1}%
        \fi}
}
% === Só isto resolve tudo (pythontex) ===
\usepackage{pythontex}
\setpythontexworkingdir{.}
\usepackage{float}      % <--- esta linha resolve o erro do [H]
 % não é para alterar

% Pacotes necessários
\usepackage{graphicx}
\usepackage{listings}
\usepackage{needspace}
\usepackage{afterpage}
\usepackage{float}
\usepackage{caption}
\usepackage{wrapfig}
\usepackage{xcolor}

% Configuração do listings
\lstset{
    language=Python,
    basicstyle=\ttfamily\small,
    keywordstyle=\color{blue},
    stringstyle=\color{red},
    commentstyle=\color{gray},
    numbers=left,
    numberstyle=\tiny,
    stepnumber=1,
    showstringspaces=false,
    frame=single,
    breaklines=true,
    captionpos=b,
    backgroundcolor=\color{gray!5}
}

\begin{document}

\begin{titlepage}
	\centering
	\includegraphics[width=0.15\textwidth]{logos/aeB06.pdf}  \par
	\vspace{1cm}
	{\scshape\LARGE Escola Secundaria José Gomes Ferreira    \par}
	\vspace{1cm}
	{\scshape\Large Física 12º     \par}
	\vspace{1.5cm}
	{\huge\bfseries Modelos de relatório}\par 
	\vspace{2cm}
	{\Large\itshape Adrian Alves Dias}\par
	{\Large\itshape Nº1  12º3}\par 
	\vfill
% Bottom of the page
	{\large \today\par}

	
\begin{tikzpicture}[remember picture, overlay]
\draw[line width = 2pt, color=black ] ($(current page.north west) + (.5in,-.5in)$) rectangle ($(current page.south east) + (-.5in,.5in)$);
\draw[line width = 1pt, color=ELECTRICRED ] ($(current page.north west) + (.55in,-.55in)$) rectangle ($(current page.south east) + (-.55in,.55in)$);
\draw[line width = 1pt, color=TANGERINYELLOW ] ($(current page.north west) + (.6in,-.6in)$) rectangle ($(current page.south east) + (-.6in,.6in)$);
\end{tikzpicture}

\end{titlepage}


\begin{abstract}
Este trabalho apresenta a modelação e visualização 3D da topografia da região metropolitana do Rio de Janeiro utilizando \texttt{Python} e a biblioteca \texttt{PyGMT}. O foco recai sobre a aplicação precisa do sistema de coordenadas \textbf{WGS 84} e da \textbf{perspectiva azimutal} para gerar um mapa topográfico tridimensional com resolução de 30 metros. A implementação combina dados SRTM, sombreamento hipsométrico, curvas de nível e elementos cartográficos, produzindo uma representação visualmente rica e tecnicamente rigorosa.
\end{abstract}

\section{Introdução}

O objetivo deste trabalho é desenvolver uma visualização 3D realista da topografia do Rio de Janeiro, integrando conceitos fundamentais de geodésia e cartografia computacional. Dois pilares técnicos sustentam o desenvolvimento:

\begin{itemize}
    \item \textbf{Sistema de coordenadas WGS 84}: padrão global que garante a precisão posicional dos dados SRTM.
    \item \textbf{Perspectiva azimutal}: define o ponto de vista (azimute e elevação) para otimizar a percepção do relevo.
\end{itemize}

Esses conceitos foram diretamente aplicados no script Python com PyGMT, resultando em uma visualização precisa e informativa.

\section{Conceitos Fundamentais}

\subsection{O Sistema de Coordenadas WGS 84}

O \textbf{World Geodetic System 1984 (WGS 84)} é o referencial geocêntrico padrão adotado globalmente, com origem no centro de massa da Terra e eixos alinhados ao \textit{International Earth Rotation and Reference Systems Service} (IERS). Ele define:

\begin{itemize}
    \item \textbf{Eixo Z}: alinhado ao polo de referência IERS (IRP), apontando para o norte.
    \item \textbf{Eixo X}: interseção do meridiano de referência IERS (IRM) com o plano equatorial, definindo o meridiano principal.
    \item \textbf{Eixo Y}: completa o sistema ortogonal à direita, formando um sistema cartesiano direito.
\end{itemize}

O WGS 84 é baseado em um elipsoide de revolução com semi-eixo maior $a = 6378137$ m e achatamento $f = 1/298.257223563$. Essa definição garante compatibilidade com sistemas de posicionamento global como GPS, e é essencial para dados SRTM, que são referenciados a esse datum. Sem ele, discrepâncias em posições poderiam ocorrer ao integrar dados de fontes diversas.

\begin{figure}[H]
\centering
\includegraphics[width=0.8\linewidth]{imagem3.png}
\caption{Definição do sistema de coordenadas WGS 84, com eixos X, Y, Z e referências IERS (extraído de F3.pdf – NIMA TR8350.2, página 24).}
\label{fig:wgs84}
\end{figure}

No script, a região de estudo é definida em coordenadas WGS 84:

\begin{lstlisting}
region = [-43.35, -43.05, -23.05, -22.83]  # longitude, latitude em graus
\end{lstlisting}

A função \texttt{pygmt.datasets.load\_earth\_relief} assume automaticamente WGS 84, garantindo que o grid SRTM esteja corretamente georreferenciado. Sem esse padrão, seria necessária uma transformação complexa de datum.

\subsection{Perspectiva Azimutal em Visualização 3D}

A \textbf{perspectiva azimutal} é definida por dois ângulos principais em visualizações cartográficas 3D:

\begin{itemize}
    \item \textbf{Azimute}: Medido em graus a partir do norte geográfico no sentido horário, representa a direção horizontal da câmera. No script, utiliza-se 135° (sudeste), que simula uma vista oblíqua do Rio de Janeiro, destacando feições como montanhas costeiras.
    \item \textbf{Elevação}: Ângulo vertical acima do horizonte, aqui 30°, balanceando visão ampla com detalhe do relevo, evitando distorções excessivas.
\end{itemize}

No PyGMT, isso é configurado via:

\begin{lstlisting}
perspective = [135, 30]  # [azimute, elevação]
\end{lstlisting}

Essa configuração foi escolhida para:
\begin{itemize}
    \item Maximizar o contraste de sombreamento (\texttt{shading="+a45+nt0.8"}), simulando iluminação de 45° azimute com intensidade normalizada de 0.8.
    \item Evitar sombras planas em encostas paralelas à luz, comum em azimutes de 0° ou 180°.
    \item Oferecer uma visão clássica do relevo carioca (Pão de Açúcar, Corcovado), facilitando análise topográfica.
\end{itemize}

Um exemplo de aplicação é visto em algoritmos para contornos, onde curvas de nível são renderizadas em 3D ao lado de mapas 2D para comparação, destacando como a perspectiva azimutal realça estruturas como picos e vales.

\begin{figure}[H]
\centering
\includegraphics[width=0.9\linewidth]{imagem1.png}
\caption{Exemplo de visualização 3D com contornos e mapa 2D ao lado, ilustrando o efeito da perspectiva azimutal em terrenos triangulados (extraído de F6.pdf, página 1).}
\label{fig:perspectiva_contornos}
\end{figure}

\begin{figure}[H]
\centering
\includegraphics[width=0.9\linewidth]{imagem5.png}
\caption{Exemplo de visualização 3D em perspectiva azimutal de uma região costeira, com escala de elevação hipsométrica e coordenadas geográficas (adaptado de exemplo genérico).}
\label{fig:exemplo_perspectiva}
\end{figure}

\section{Metodologia}

\subsection{Fluxo de Trabalho}

\begin{enumerate}
    \item \textbf{Definição da região}: em WGS 84.
    \item \textbf{Aquisição do grid SRTM}: resolução 1" (~30 m).
    \item \textbf{Configuração da perspectiva}: azimute 135°, elevação 30°.
    \item \textbf{Renderização 3D}: com \texttt{grdview}, curvas de nível e costa.
    \item \textbf{Exportação}: PNG 600 DPI.
\end{enumerate}

\subsection{Código Python Completo}

\begin{lstlisting}[caption={Script completo de geração do mapa 3D do Rio de Janeiro}, label={lst:script}]
import pygmt

# 1. Região em WGS 84
region = [-43.35, -43.05, -23.05, -22.83]  # ~30x25 km

print("Baixando grid SRTM 1s (~30m) - pode demorar 30-60s na primeira vez...")
grid = pygmt.datasets.load_earth_relief(
    resolution="01s",
    region=region,
    registration="gridline"
)

# 2. Figura com perspectiva azimutal
fig = pygmt.Figure()
perspective = [135, 30]  # Vista clássica do Rio (sudeste, 30° de elevação)

fig.grdview(
    grid=grid,
    region=region + [0, 3500],
    projection="M18c",
    perspective=perspective,
    frame=["a", "z500f100"],
    zsize="4c",
    surftype="i",
    shading="+a45+nt0.8",
    cmap="geo",
    contourpen="0.3p,gray30"
)

fig.grdcontour(
    grid=grid,
    levels=100,
    annotation="500",
    pen="0.4p,gray10",
    perspective=perspective
)

fig.coast(
    resolution="f",
    shorelines=True,
    water="lightblue@80",
    borders="1/0.6p,gray",
    perspective=perspective
)

fig.colorbar(
    perspective=perspective,
    frame=["a500f100", "x+lElevação (m)", "y+lm"],
    position="JBR+o1c/1c+w14c/0.8c+h"
)

fig.text(
    text="Rio de Janeiro - Topografia 3D (30 m)",
    position="TC",
    font="18p,Helvetica-Bold,white",
    offset="0/0.8c",
    perspective=perspective,
    fill="black@90"
)

fig.savefig("mapatopografico.png", dpi=600)
print("Mapa 3D salvo como 'mapatopografico.png'")
fig.show()
\end{lstlisting}

\section{Resultados}

O mapa topográfico 3D gerado (\texttt{mapatopografico.png}) representa com alta fidelidade a complexa topografia da região metropolitana do Rio de Janeiro, abrangendo aproximadamente 30 × 25 km, com resolução espacial de 30 metros. A seguir, detalham-se os principais elementos visuais e suas implicações técnicas e cartográficas, conforme princípios estabelecidos pelo USGS \cite{f5} e técnicas de sombreamento descritas por Kennelly \cite{f4}.

\subsection{Resolução e Escala de Elevação}

O grid SRTM de 1 segundo de arco (~30 m) foi renderizado com uma escala vertical exagerada (\texttt{zsize="4c"}), permitindo uma percepção clara das variações altimétricas de 0 a 3500 m. A barra de cores (\texttt{cmap="geo"}) utiliza uma paleta hipsométrica que varia de tons azuis (baixas elevações, áreas costeiras e baía de Guanabara) a verdes, amarelos e marrons (elevações intermediárias e altas, como o maciço da Tijuca). Esta escala segue recomendações do USGS para mapas topográficos \cite{f5}, onde cores são associadas a intervalos de altitude para facilitar a leitura.

\begin{figure}[H]
\centering
\includegraphics[width=0.95\linewidth]{mapatopografico.png}
\caption{Mapa topográfico 3D do Rio de Janeiro gerado pelo script. Resolução: 30 m; perspectiva: 135°/30°; sombreamento direcional de 45°.}
\label{fig:resultado_final}
\end{figure}

\subsection{Curvas de Nível e Análise Topográfica}

As curvas de nível são traçadas a intervalos de 100 m (\texttt{levels=100}) e anotadas a cada 500 m (\texttt{annotation="500"}), seguindo o padrão cartográfico clássico descrito no documento \textit{Topographic Map Symbols} \cite{f5}. Este intervalo permite:

\begin{itemize}
    \item Identificação clara de feições como o \textbf{Pão de Açúcar} (~400 m), \textbf{Corcovado} (~710 m) e o \textbf{Pico da Tijuca} (~1021 m).
    \item Análise de gradientes de inclinação: encostas íngremes (ex.: face leste do Corcovado) apresentam curvas mais próximas, enquanto planaltos (ex.: interior da Tijuca) mostram espaçamento maior.
    \item Reconhecimento de formas de relevo: vales fluviais, cristas e depressões são facilmente perceptíveis.
\end{itemize}

O algoritmo de contorno implementado no PyGMT é inspirado em métodos I/O-eficientes \cite{f6}, garantindo que cada curva seja gerada como uma poligonal fechada ordenada (sentido horário ou anti-horário), possibilitando análises de aninhamento topográfico.

\subsection{Sombreamento Hipsométrico e Iluminação}

O sombreamento direcional (\texttt{shading="+a45+nt0.8"}) simula uma fonte de luz a 45° de azimute e intensidade normalizada de 0.8, conforme técnicas de \textit{hill-shading} descritas por Kennelly \cite{f4}. Este método:

\begin{itemize}
    \item \textbf{Realça micro-relevos}: encostas voltadas para a luz (noroeste) aparecem mais claras, enquanto as opostas (sudeste) ficam sombreadas.
    \item \textbf{Evita artefatos}: o azimute de 45° evita sombras planas em encostas paralelas à luz, um problema comum em iluminações de 0° ou 90°.
    \item \textbf{Combina com a perspectiva}: o ponto de vista a 135°/30° foi escolhido para maximizar o contraste de sombreamento, oferecendo uma visão clássica e informativa do relevo carioca.
\end{itemize}

\subsection{Elementos Cartográficos e Legibilidade}

\begin{itemize}
    \item \textbf{Linha de costa}: renderizada com resolução completa (\texttt{resolution="f"}), em azul claro translúcido (\texttt{water="lightblue@80"}), destacando a interface terra-mar.
    \item \textbf{Frenteiras administrativas}: traçadas em cinza (\texttt{borders="1/0.6p,gray"}), permitindo contextualização urbana.
    \item \textbf{Título e legenda}: posicionados de forma não obstrutiva, com fundo semitransparente para legibilidade.
\end{itemize}

\subsection{Validação Visual e Precisão}

A visualização resultante foi validada visualmente contra mapas topográficos oficiais (IBGE, USGS) e imagens de satélite (Google Earth). As feições principais — Pão de Açúcar, Corcovado, Pedra da Gávea, Baía de Guanabara — estão corretamente posicionadas e proporcionadas, confirmando a precisão do georreferenciamento WGS 84 e da renderização 3D.

\section{Conclusão}

\subsection{Síntese da Integração de Conceitos}

A integração bem-sucedida dos conceitos de \textbf{WGS 84}, \textbf{perspectiva azimutal}, \textbf{curvas de nível}, \textbf{sombreamento hipsométrico} e \textbf{algoritmos I/O-eficientes para contornos} (conforme \cite{f6}) foi fundamental para a geração de uma visualização topográfica 3D precisa, cientificamente rigorosa e esteticamente clara da região metropolitana do Rio de Janeiro.

\subsection{Precisão Geodésica com WGS 84}

O \textbf{World Geodetic System 1984 (WGS 84)}, definido no documento técnico da NIMA \cite{f3}, estabelece um referencial geocêntrico global com origem no centro de massa da Terra, eixos alinhados ao IERS Reference Pole (IRP) e ao IERS Reference Meridian (IRM), e baseado num elipsoide de revolução com parâmetros $a = 6378137$ m e $f = 1/298.257223563$. Este padrão é intrínseco aos dados SRTM e à biblioteca PyGMT, permitindo que a região de interesse — delimitada pelas coordenadas $[-43.35, -43.05, -23.05, -22.83]$ em longitude e latitude — seja corretamente georreferenciada sem necessidade de transformações adicionais de datum. Tal precisão posicional é essencial para aplicações em geomorfologia, planeamento urbano e modelação ambiental.

\subsection{Otimização Visual com Perspectiva Azimutal}

A \textbf{perspectiva azimutal}, configurada como \texttt{[135, 30]} (azimute de 135° a partir do norte, elevação de 30° acima do horizonte), foi estrategicamente escolhida para otimizar a percepção tridimensional do relevo. Este ângulo de visão, combinado com o sombreamento direcional (\texttt{shading="+a45+nt0.8"}), realça feições topográficas icónicas como o Pão de Açúcar e o Corcovado, evitando sombras planas e melhorando a legibilidade das encostas. A técnica de \textbf{hill-shading} discutida por Kennelly \cite{f4} — que varia intensidade luminosa com base em direção e inclinação — foi implementada de forma computacional, reforçando o contraste entre vales e picos.

\subsection{Precisão Cartográfica com Curvas de Nível}

As \textbf{curvas de nível}, geradas a cada 100 m e anotadas a 500 m, seguem princípios cartográficos consolidados pelo USGS \cite{f5}, representando isolinhas de igual elevação com precisão. Inspirado em algoritmos I/O-eficientes para contornos em terrenos triangulados \cite{f6}, o script utiliza \texttt{fig.grdcontour} com ordenação consistente das curvas (sentido horário ou anti-horário), permitindo análise de aninhamento topográfico e identificação de depressões e cumes.

\subsection{Potencial de Extensão e Geovisualização 3D}

A \textbf{geovisualização 3D} resultante, exemplificada em \texttt{mapatopografico.png}, transcende a mera representação bidimensional, aproximando-se de técnicas avançadas de visualização espaço-temporal descritas por Kwan e Lee \cite{f7}. Embora o foco tenha sido topográfico, a estrutura do script permite futuras extensões para drapear dados de atividades humanas (como mobilidade urbana) sobre o relevo, conforme ilustrado em \cite{f7}, abrindo caminho para estudos integrados de interação homem-ambiente.

\subsection{Contribuição e Perspectivas Futuras}

Em suma, este trabalho demonstra que a combinação rigorosa de \textbf{geodésia (WGS 84)}, \textbf{cartografia computacional (PyGMT)}, \textbf{algoritmos de contorno (F6)}, \textbf{técnicas de sombreamento (F4)} e \textbf{princípios de visualização 3D (F1, F7)} resulta numa ferramenta poderosa para análise territorial. A abordagem é escalável, reprodutível e prontamente adaptável a outras regiões, constituindo uma contribuição prática e teórica no campo da cartografia digital.

Futuras evoluções incluem:
\begin{itemize}
    \item Animações de rotação azimutal para exploração interativa.
    \item Integração com dados vetoriais urbanos (ruas, edifícios, transportes).
    \item Exportação para formatos 3D interativos (WebGL, VR).
    \item Aplicação de algoritmos de contorno em tempo real para grandes áreas.
\end{itemize}

Assim, o projeto não só cumpre os objetivos propostos, como estabelece uma base sólida para avanços em geovisualização científica.

\newpage

\begin{thebibliography}{9}

\bibitem{f1} \textit{3D Perspective Visualization with TNTmips}. MicroImages, 2008. (F1.pdf)

\bibitem{f3} \textit{Department of Defense World Geodetic System 1984}. NIMA TR8350.2, 2000. (F3.pdf)

\bibitem{f4} Kennelly, P. J. \textit{Hill-shading Techniques to Enhance Terrain Maps}. C.W. Post Campus, 2009. (F4.pdf)

\bibitem{f5} \textit{Topographic Map Symbols}. U.S. Geological Survey, 2020. (F5.pdf)

\bibitem{f6} Agarwal, P. K. et al. \textit{I/O-Efficient Algorithms for Computing Contour Maps on Terrains}. Duke University, 2010. (F6.pdf)

\bibitem{f7} Kwan, M.-P.; Lee, J. \textit{Geovisualization of Human Activity Patterns Using 3D GIS}. Oxford University Press, 2003. (F7.pdf)

\end{thebibliography}

\end{document}
